\chapter{प्रत्यावर्ती धारा}\label{ch:g9:alternating-current}

तुमने अब तक जो भी परिपथ बनाया, वह \omterm{def:g3:first-electric-circuit:battery}{बैटरी} पर चला: धारा टिककर $+$ से $-$
की ओर क़तार बाँधे चलती हुई, घड़ी की सुई जितनी फ़र्ज़-निभाऊ। पर तुम्हारी
दीवारों के सॉकेट कुछ और ही जंगली चीज़ पहुँचाते हैं --- एक ऐसी धारा जो
हर सेकंड सौ बार अपनी दिशा उलट देती है, और जिसका \omterm{def:g8:voltage-voltmeter:voltage}{विभवांतर} बिना रुके
धन और ऋण के बीच झूलता रहता है। मिलो \omterm{def:g9:alternating-current:dcac}{प्रत्यावर्ती धारा} से: वही रूप,
जिसमें दुनिया की लगभग सारी बिजली जन्मती है, ढोई जाती है और पहुँचाई
जाती है।

\section{धारा के दो क़िस्म}

\begin{definition}[दिष्ट और प्रत्यावर्ती धारा]\label{def:g9:alternating-current:dcac}
\emph{दिष्ट धारा}\index{दिष्ट धारा} (DC) हमेशा एक ही दिशा में बहती है,
और उसे टिका हुआ \omterm{def:g8:voltage-voltmeter:voltage}{विभवांतर} चलाता है --- यानी \omterm{def:g3:first-electric-circuit:battery}{बैटरी} की देन, जो भंडार
चुकने तक अटल रहती है। जबकि \emph{प्रत्यावर्ती धारा}\index{प्रत्यावर्ती
धारा} (AC) समय-समय पर अपनी दिशा उलट देती है, और उसे वह \omterm{def:g8:voltage-voltmeter:voltage}{विभवांतर} चलाता
है जो धन से ऋण और वापस, बार-बार, मुलायम ढंग से झूलता रहता है ---
सॉकेट का राज, और प्रत्यावर्तित्र का भी (किसका, यह अगला अध्याय बताता
है)।
\end{definition}

\begin{omfigure}
\begin{tikzpicture}
\begin{axis}[omaxis, width=10.8cm, height=5.2cm,
    xlabel={समय (\unit{ms})}, ylabel={विभवांतर (\unit{V})},
    xtick={0,10,20,30,40}, ytick={-325,0,230,325},
    yticklabels={$-325$,$0$,$230$,$325$},
    ymin=-400, ymax=400, xmin=0, xmax=42, samples=200]
  \addplot[very thick, omDef, domain=0:40] {230*0+4.5};
  \addplot[very thick, omThm, domain=0:40]
    {325*sin(deg(2*3.14159*0.05*x))};
  \node[right, font=\footnotesize, omDef] at (axis cs:29,42)
    {बैटरी: टिकी हुई \qty{4.5}{V}};
  \node[right, font=\footnotesize, omThm] at (axis cs:1.2,352)
    {मेन: झूलता हुआ};
\end{axis}
\end{tikzpicture}

{\small समय के सामने दो \omterm{def:g8:voltage-voltmeter:voltage}{विभवांतर}: \omterm{def:g3:first-electric-circuit:battery}{बैटरी} की सपाट, वफ़ादार रेखा --- और
मेन का धन तथा ऋण के बीच अनंत झूला, हर पूरे चक्र पर बीस मिलीसेकंड।}
\end{omfigure}

\begin{definition}[आवर्तकाल और आवृत्ति]\label{def:g9:alternating-current:period}
किसी प्रत्यावर्ती \omterm{def:g8:voltage-voltmeter:voltage}{विभवांतर} का \emph{आवर्तकाल}\index{आवर्तकाल} $T$ एक
पूरे चक्र की \omterm{def:g2:measuring-time:duration}{अवधि} है --- यानी एक पूरा झूला ऊपर, नीचे और वापस, सेकंडों
में। और उसकी \emph{\omterm{def:g8:sound-pitch-loudness:frequency}{आवृत्ति}} $f$ हर सेकंड के चक्रों की गिनती है,
\omterm{def:g8:sound-pitch-loudness:frequency}{हर्ट्ज़} में --- \omterm{def:g3:sound-around-us:sound}{ध्वनि} का वही \omterm{def:g6:measurement-in-science:measuring}{मात्रक}, फिर से सेवा में --- और ये
दोनों एक-दूसरे के व्युत्क्रम हैं:
\[
f = \frac{1}{T}, \qquad T = \frac{1}{f}.
\]
\end{definition}

\begin{proposition}[मेन के ज़रूरी आँकड़े]\label{prop:g9:alternating-current:mains}
दुनिया के ज़्यादातर हिस्सों में मेन $f = \qty{50}{Hz}$ पर प्रत्यावर्तन
करता है --- यानी हर सेकंड पचास पूरे चक्र, और \omterm{def:g9:alternating-current:period}{आवर्तकाल} $T = 1/50 =
\qty{0.02}{s}$: हर झूले पर बीस मिलीसेकंड (कुछ इलाक़े साठ पर चलते हैं)।
तुम्हारे पढ़ने वाले लैंप की डोरी में धारा हर सेकंड सौ बार अपनी दिशा
उलटती है --- हर चक्र पर दो बार --- और उसने एक बार भी एक सिरे से
दूसरे सिरे तक टिककर क़तार नहीं बाँधी।
\end{proposition}

\begin{example}[आवृत्ति का अंकगणित]\label{ex:g9:alternating-current:arithmetic}
यह व्युत्क्रम वाली जोड़ी दोनों ओर हिसाब लगा देती है। \qty{60}{Hz} पर
चक्कर काटता कोई जनित्र: $T = 1/60 \approx \qty{0.017}{s}$ --- यानी
सत्रह मिलीसेकंड। कोई \omterm{met:g9:alternating-current:scope}{दोलनदर्शी} \qty{0.04}{s} लंबा एक चक्र दिखाता है:
$f = 1/0.04 = \qty{25}{Hz}$. धीमे पैडल पर साइकिल का डाइनेमो
\qty{10}{Hz} बनाता है: यानी \omterm{def:g9:alternating-current:period}{आवर्तकाल} सेकंड का दसवाँ हिस्सा --- और उसका
लैंप साफ़-साफ़ टिमटिमाता है, हर सेकंड दस स्पंद, जब तक रफ़्तार उन्हें
आँख की पकड़ से परे जाकर चिकना न कर दे।
\end{example}

\section{झूले को पढ़ना}

\begin{method}[दोलनदर्शी की रेखा पढ़ना]\label{met:g9:alternating-current:scope}
\emph{दोलनदर्शी}\index{दोलनदर्शी} --- या कोई भी जाँचने वाला अनुप्रयोग
--- खानों वाले पर्दे पर \omterm{def:g8:voltage-voltmeter:voltage}{विभवांतर} को समय के सामने खींचता है; और हर खाना
कितने \omterm{def:g8:voltage-voltmeter:voltage}{वोल्ट} तथा कितने मिलीसेकंड का है, यह दो घुंडियाँ तय करती हैं। कोई
रेखा काटने के लिए:
\begin{enumerate}
  \item एक पूरे चक्र के खाने गिनो और प्रति-खाना समय से गुणा करो: वही
        $T$ है; और फिर $f = 1/T$;
  \item बीच की रेखा से किसी शिखर तक के खाने गिनो और प्रति-खाना \omterm{def:g8:voltage-voltmeter:voltage}{वोल्ट}
        से गुणा करो: यही \emph{शिखर \omterm{def:g8:voltage-voltmeter:voltage}{विभवांतर}} $U_{\max}$ है;
  \item आकृति जाँचो: मेन और प्रत्यावर्तित्र चित्र वाली वही मुलायम,
        अनंत बार दोहराई जाने वाली तरंग खींचते हैं ---
        \emph{ज्यावक्र}, जो उस पूर्ण झूले को गणित का दिया नाम है।
\end{enumerate}
\end{method}

\begin{proposition}[प्रभावी विभवांतर]\label{prop:g9:alternating-current:effective}
जो झूला हर सेकंड सौ बार $+U_{\max}$ से $-U_{\max}$ तक के
हर मान से होकर गुज़रता हो, उसका ``वह'' \omterm{def:g8:voltage-voltmeter:voltage}{विभवांतर} आख़िर है क्या? ईमानदार
अकेला अंक \emph{प्रभावी विभवांतर}\index{प्रभावी विभवांतर} $U$ है:
यानी वह टिका हुआ DC \omterm{def:g8:voltage-voltmeter:voltage}{विभवांतर}, जो किसी \omterm{def:g8:ohms-law:resistor}{प्रतिरोधक} को उतना ही गरम करता।
ज्यावक्र के लिए वह शिखर से काफ़ी नीचे बैठता है --- $U_{\max}$
$U$ का क़रीब $1.4$ गुना होता है --- और हर AC \omterm{def:g8:voltage-voltmeter:voltmeter}{वोल्टमीटर} इसी को पढ़ने के
लिए बना होता है। सॉकेट की वह मशहूर ``\qty{230}{V}'' प्रभावी मान है;
जबकि शिखर सचमुच क़रीब \qty{325}{V} तक पहुँचते हैं, हर सेकंड पचास बार,
दोनों ओर।
\end{proposition}

\begin{example}[किन उपकरणों को फ़र्क़ पड़ता है]\label{ex:g9:alternating-current:whocares}
\omterm{prop:g8:electric-current-effects:three}{ऊष्मीय प्रभाव} दोनों ही ढंग से ठीक काम करता है --- कोई तंतु या केतली की
कुंडली क़तार की दिशा की परवाह किए बिना गरम होती है, और इसीलिए लैंप तथा
हीटर बिना शिकायत AC पी लेते हैं। पर इलेक्ट्रॉनिकी --- फ़ोन, लैपटॉप,
चिप वाली हर चीज़ --- को टिकी हुई DC चाहिए: उनकी चार्जर वाली ईंटें और
छिपी हुई विद्युत आपूर्तियाँ सॉकेट के झूले को \omterm{def:g3:first-electric-circuit:battery}{बैटरी} जैसी अटलता में बदल
देती हैं (ईंट की गरमाहट टटोलो: बदलाव अपनी चुंगी वसूल लेता है)। और
बैटरियों की दुनिया भी झूले पर चार्ज नहीं हो सकती: चार्ज होना उलटी दिशा
में चलता रसायन है, और रसायन एक टिकी हुई दिशा पर अड़ जाता है।
\end{example}

\begin{remark}[दुनिया ने झूला क्यों चुना]\label{rem:g9:alternating-current:why}
अगर बैटरियाँ इतनी सुखद ढंग से टिकी हुई हैं, तो सभ्यता की तारबंदी AC के
लिए क्यों है? दो वजहें, और दोनों दो अध्याय आगे इंतज़ार कर रही हैं:
\omterm{def:g9:alternating-current:dcac}{प्रत्यावर्ती धारा} वही है जो घूमती मशीनें \emph{क़ुदरतन बनाती} हैं ---
और, सबसे बढ़कर, ट्रांसफ़ॉर्मर को सिर्फ़ AC ही खिलाती है, यानी वह
मामूली-सा गुनगुनाता डिब्बा जो बिजली को विशाल \omterm{def:g8:voltage-voltmeter:voltage}{विभवांतर} पर पूरा देश पार
करने देता है और तुम्हारे घर में सधे हुए \omterm{def:g8:voltage-voltmeter:voltage}{विभवांतर} पर घुसने देता है। यह
जाल इसलिए प्रत्यावर्ती है कि सफ़र यही माँगता है; और अगले दो अध्याय उस
धारा के पीछे-पीछे घूमती कुंडली से रसोई की दीवार तक चलते हैं।
\end{remark}

\begin{remark}[झूला और शरीर]\label{rem:g9:alternating-current:safety}
सुरक्षा के पुराने अंकगणित को एक भयावह ब्योरा और मिल जाता है: मेन के
हर सेकंड के पचास झूले उस लय के पास पड़ते हैं जो दिल की अपनी विद्युत
समय-सारणी सबसे आसानी से बिगाड़ देती है --- यानी \omterm{def:g8:intensity-ammeter:intensity}{ऐम्पियर} दर \omterm{def:g8:intensity-ammeter:intensity}{ऐम्पियर},
प्रत्यावर्ती झटके टिके हुए झटकों से ज़्यादा ख़तरनाक होते हैं। बचपन का
हर नियम अपनी जगह क़ायम है, और दोगुना होकर: सॉकेट की प्रभावी
\qty{230}{V} --- शिखर पर \qty{325}{V} --- सिर्फ़ कोई बड़ी \omterm{def:g3:first-electric-circuit:battery}{बैटरी}
नहीं है। वह एक अलग, और कम माफ़ करने वाला जानवर है।
\end{remark}

\section{अभ्यास}

\begin{exercise}[$\star$]\label{exo:g9:alternating-current:1}
DC और AC में एक-एक वाक्य में फ़र्क़ करो, और हर एक का घरेलू स्रोत भी
बताओ।
\end{exercise}

\begin{exercise}[$\star$]\label{exo:g9:alternating-current:2}
\omterm{def:g9:alternating-current:period}{आवर्तकाल} और \omterm{def:g8:sound-pitch-loudness:frequency}{आवृत्ति} की परिभाषा दो और उनका व्युत्क्रम वाला संबंध लिखो।
\qty{50}{Hz} के लिए $T$ निकालो, और $T = \qty{0.01}{s}$ के लिए $f$.
\end{exercise}

\begin{exercise}[$\star$]\label{exo:g9:alternating-current:3}
\qty{50}{Hz} पर मेन की धारा हर सेकंड कितनी बार अपनी दिशा उलटती है?
(सावधान: हर चक्र पर कितने उलटाव।)
\end{exercise}

\begin{exercise}[$\star$]\label{exo:g9:alternating-current:4}
कोई \omterm{met:g9:alternating-current:scope}{दोलनदर्शी} \qty{5}{ms} प्रति खाना पर एक चक्र $4$ खानों में फैला
दिखाता है। $T$ और $f$ निकालो।
\end{exercise}

\begin{exercise}[$\star$]\label{exo:g9:alternating-current:5}
सॉकेट \qty{230}{V} कहता है; और शिखर \qty{325}{V} तक पहुँचता है। दोनों
मानों के नाम बताओ, और छोटे वाले का मतलब भी।
\end{exercise}

\begin{exercise}[$\star$]\label{exo:g9:alternating-current:6}
केतलियाँ AC ख़ुशी से क़बूल कर लेती हैं जबकि फ़ोन के चार्जरों को उसे
बदलना पड़ता है --- क्यों? एक प्रभाव, और एक रसायन।
\end{exercise}

\begin{exercise}[$\star$]\label{exo:g9:alternating-current:7}
दोनों रेखाएँ खींचो (या बयान करो): \qty{9}{V} की \omterm{def:g3:first-electric-circuit:battery}{बैटरी}, और
\qty{20}{ms} \omterm{def:g9:alternating-current:period}{आवर्तकाल} तथा \qty{9}{V} शिखर वाला प्रत्यावर्ती \omterm{def:g8:voltage-voltmeter:voltage}{विभवांतर}।
उपकरण की खींची कौन-सी अकेली ख़ासियत उन्हें एक नज़र में अलग बता देती
है?
\end{exercise}

\begin{exercise}[$\star$]\label{exo:g9:alternating-current:8}
अध्याय की झलक के मुताबिक़ सभ्यता ने अपनी तारबंदी AC के लिए क्यों की ---
दोनों वजहें बताओ।
\end{exercise}

\begin{exercise}[$\star\star$]\label{exo:g9:alternating-current:9}
धीमे चलता साइकिल का डाइनेमो अपना लैंप साफ़ टिमटिमाहट के साथ जलाता है;
और ज़ोर से पैडल मारने पर दमक टिक जाती है। दोनों हालतें \omterm{def:g8:sound-pitch-loudness:frequency}{आवृत्ति} से
समझाओ --- और यह भी अंदाज़ो कि किस \omterm{def:g8:sound-pitch-loudness:frequency}{आवृत्ति} से नीचे टिमटिमाहट खलने लगती
है, यह जानते हुए कि सिनेमा क़रीब $24$ तसवीरें प्रति सेकंड पर आँख को
बहला लेता है।
\end{exercise}

\begin{exercise}[$\star\star$]\label{exo:g9:alternating-current:10}
किसी सॉकेट पर लगा AC \omterm{def:g8:voltage-voltmeter:voltmeter}{वोल्टमीटर} \qty{230}{V} बताता है। उसी सॉकेट पर लगा
कोई (काल्पनिक) निडर \omterm{met:g9:alternating-current:scope}{दोलनदर्शी} शिखर किस मान पर दिखाएगा --- और दोनों
चिह्न गिनते हुए हर सेकंड कितने शिखर?
\end{exercise}

\begin{exercise}[$\star\star$]\label{exo:g9:alternating-current:11}
किसी इलाक़े का जाल \qty{60}{Hz} पर, प्रभावी \qty{120}{V} के साथ चलता
है। उसका \omterm{def:g9:alternating-current:period}{आवर्तकाल} और लगभग शिखर \omterm{def:g8:voltage-voltmeter:voltage}{विभवांतर} निकालो --- और यह भी बताओ कि
तुम्हारे सफ़र के किन उपकरणों को एतराज़ होगा
(\cref{ex:g9:alternating-current:whocares} देखो; आधुनिक चार्जर अपने
लेबलों की बारीक इबारत पढ़ लेते हैं)।
\end{exercise}

\begin{exercise}[$\star\star\star$]\label{exo:g9:alternating-current:12}
पर्दों और फ़्लोरोसेंट लैंपों की पुरानी फ़िल्मों में कभी-कभी काली
लुढ़कती पट्टियाँ दिखती हैं --- यानी कैमरे के फ़्रेम लैंप की चमक को झूले
के बीचोंबीच पकड़ लेते हैं। \qty{50}{Hz} पर चलता मेन का लैंप
\emph{दोनों} चिह्नों के हर शिखर पर चमक उठता है: तो हर सेकंड कितने
चमक-स्पंद? और $25$ फ़्रेम प्रति सेकंड वाला कैमरा उस स्पंदन का नमूना
लेता है: तर्क करो कि दोनों लयों के बीच की टकराहट धीरे-धीरे सरकती
पट्टियाँ क्यों पोत देती है, और \qty{60}{Hz} वाले इलाक़े में फ़िल्म
बनाने से वह नमूना क्यों बदल जाता है।
\end{exercise}

\section{समस्या: दोलनदर्शी की मेज़}

\begin{problem}[{सप्ताहांत समस्या --- दोलनदर्शी की मेज़ पर योग्यता की
रात; चार पहेली-आपूर्तियाँ, एक पर्दा; तकनीशियन का
बहीखाता}]\label{pb:g9:alternating-current:1}
बिना लेबल वाली चार आपूर्तियों की पहचान उनकी रेखाओं से करनी है।
\omterm{met:g9:alternating-current:scope}{दोलनदर्शी} की घुंडियाँ: आड़े हर खाने पर \qty{5}{ms}, और खड़े हर खाने पर
\qty{2}{V} (जहाँ अलग से न कहा गया हो)। बहीखाता तुम रखते हो।

\medskip\noindent\textbf{भाग I --- चारों रेखाएँ।}
\begin{enumerate}
  \item आपूर्ति A: एक सपाट रेखा, बीच से $2.25$ खाने ऊपर। उसका क़िस्म
        और \omterm{def:g8:voltage-voltmeter:voltage}{विभवांतर} बताओ --- और उसकी संभावित घरेलू पहचान भी।
  \item आपूर्ति B: एक मुलायम तरंग, हर $4$ खानों में एक चक्र, और शिखर
        $3$ खाने ऊँचे। क़िस्म, \omterm{def:g9:alternating-current:period}{आवर्तकाल}, \omterm{def:g8:sound-pitch-loudness:frequency}{आवृत्ति}, शिखर \omterm{def:g8:voltage-voltmeter:voltage}{विभवांतर}?
  \item आपूर्ति C: बीच से $2.25$ खाने \emph{नीचे} एक सपाट रेखा। यह
        क्या है, और कौन-सी नन्ही ग़लती इसे समझा देती है?
  \item आपूर्ति D: एक मुलायम तरंग, हर $2$ खानों में एक चक्र, और शिखर
        $1.5$ खाने। पूरे आँकड़े --- और यह भी कि अगर दोनों एक-एक लैंप
        खिलाएँ तो D का लैंप B के लैंप से कम क्यों टिमटिमाता।
\end{enumerate}

\medskip\noindent\textbf{भाग II --- तकनीशियन के सवाल।}
\begin{enumerate}[resume]
  \item आपूर्ति B पर लगा AC \omterm{def:g8:voltage-voltmeter:voltmeter}{वोल्टमीटर} क़रीब \qty{4.2}{V} बताता है,
        जबकि उसके शिखर \qty{6}{V} तक पहुँचते हैं। दोनों अंकों को
        मिलाओ।
  \item चारों में से कौन-सी आपूर्ति किसी फ़ोन को सीधे चार्ज कर सकती
        है, और बाक़ियों तथा उसकी \omterm{def:g3:first-electric-circuit:battery}{बैटरी} के बीच क्या खड़ा होना चाहिए?
  \item मेज़ की चेतावनी वाली पट्टी: ``AC \omterm{def:g8:voltage-voltmeter:voltage}{वोल्ट} दर \omterm{def:g8:voltage-voltmeter:voltage}{वोल्ट} ज़्यादा
        काटती है।'' उसकी वजह दोबारा कहो।
  \item एक प्रशिक्षु पूछता है कि \omterm{met:g9:alternating-current:scope}{दोलनदर्शी} B की तरंग टेढ़ी-मेढ़ी के
        बजाय इतनी मुलायम गोलाई में क्यों खींचता है। उस आकृति का नाम
        और उसका उद्गम बताओ (अगले अध्याय की कौन-सी मशीन उसे क़ुदरतन
        खींचती है?)।
\end{enumerate}

\medskip\noindent\textbf{भाग III --- ख़ुद जाल की रेखा।}
सुरक्षित ढंग से अलग किया गया एक प्रदर्शक मेन दिखाता है, \qty{100}{V}
और \qty{5}{ms} प्रति खाना पर।
\begin{enumerate}[resume]
  \item रेखा की भविष्यवाणी करो: हर चक्र पर कितने खाने, और शिखर की
        ऊँचाई कितने खाने।
  \item उसी रेखा से $f$ और \omterm{prop:g9:alternating-current:effective}{प्रभावी विभवांतर} निकालो --- और दोनों को
        सॉकेट के मशहूर लेबल के सामने जाँचो भी।
  \item प्रदर्शक को बराबर प्रभावी ऊष्मा-शक्ति वाली बैटरियों की क़तार
        पर बदल दिया जाता है। नई रेखा का बयान एक पंक्ति में करो।
  \item बहीखाता तकनीशियन के तीन पंक्तियों वाले सार से बंद करो: सपाट
        रेखा का क्या मतलब है, तरंग का क्या, और कौन-सा अकेला अंक दोनों
        की इंसाफ़ वाली तुलना करने देता है।
\end{enumerate}
\end{problem}
